% GENERATED FILE. Edit canonical lessons, then run npm run generate:books.
% canonical-source-hash: fnv1a64:f2fba7741a357488
% canonical-lessons: TE-C75-today, TE-C75-yesterday, TE-C75-when, TE-C75-before, TE-C75-year

\chapter{Today, Yesterday, When}
\label{ch:te-today-yesterday-when}
\hlchaptermodality{75}

\begin{chapteropening}
\textbf{By the end of this chapter:} \emph{I can say today, yesterday and before in Telugu, ask when something happens, and name a year.}

Complete TE-C75-year and demonstrate the chapter promise.
\end{chapteropening}

\section[ivāḷa]{\te{ఇవాళ} (ivāḷa) — today}

\label{lesson:TE-C75-today}



\begin{quote}
Before the new one: what did \emph{endukaṇṭē} mean?
\end{quote}

\subsection*{You'll want to know: \te{ఇవాళ}}
\textbf{\te{ఇవాళ}} (\emph{ivāḷa}) — \textquotedblleft{}today\textquotedblright{}.

The farewells gave you \te{రేపు}, \textquotedblleft{}tomorrow\textquotedblright{}, inside \te{రేపు} \te{కలుద్దాం}. The day you are actually standing in had no name until now.

It begins with the same \textbf{\te{ఇ}-} as \te{ఇక్కడ} and \te{ఇది}: the near one, the one right here. \te{ఇవాళ} is the near day.

\begin{grammarlens}[title={What you've built}]
The day you are in.
\end{grammarlens}

\subsection*{Your turn}
\begin{itemize}
\raggedright
  \item \emph{Say it:} \emph{ivāḷa}
  \item \emph{Say it:} \emph{ivāḷa}, once more
  \item \emph{Say it:} say \emph{ivāḷa}, then \emph{rēpu}, and say which is nearer
  \item \emph{Recall:} say \emph{bhōjanaṁ}, then say \emph{ivāḷa}
\end{itemize}

\begin{tcolorbox}[breakable,colback=teal!4,colframe=teal!35!black,title={Before you move on}]
What does \te{ఇవాళ} mean? (\textquotedblleft{}today\textquotedblright{}.) Say it once more.
\end{tcolorbox}
\section[ninna]{\te{నిన్న} (ninna) — yesterday}

\label{lesson:TE-C75-yesterday}



\begin{quote}
Before the new one: what did \emph{ivāḷa} mean?
\end{quote}

\subsection*{You'll want to know: \te{నిన్న}}
\textbf{\te{నిన్న}} (\emph{ninna}) — \textquotedblleft{}yesterday\textquotedblright{}.

With \te{ఇవాళ} and \te{రేపు} you now have the three days every conversation stands on, and \te{నిన్న} is the only one that looks backward.

It is native Dravidian and short, with the doubled \textbf{\te{న్న}} you have seen in \te{ఉన్నాను}.

\begin{grammarlens}[title={What you've built}]
Three days: \te{నిన్న}, \te{ఇవాళ}, \te{రేపు}.
\end{grammarlens}

\subsection*{Your turn}
\begin{itemize}
\raggedright
  \item \emph{Say it:} \emph{ninna}
  \item \emph{Say it:} \emph{ninna}, once more
  \item \emph{Say it:} say the three days in order: \emph{ninna}, \emph{ivāḷa}, \emph{rēpu}
  \item \emph{Recall:} say \emph{vātāvaraṇam}, then say \emph{ninna}
\end{itemize}

\begin{tcolorbox}[breakable,colback=teal!4,colframe=teal!35!black,title={Before you move on}]
What does \te{నిన్న} mean? (\textquotedblleft{}yesterday\textquotedblright{}.) Say it once more.
\end{tcolorbox}
\section[eppuḍu vastāru?]{\te{ఎప్పుడు} \te{వస్తారు}? (eppuḍu vastāru?) — when will you come?}

\label{lesson:TE-C75-when}



\begin{quote}
Before the new one: what did \emph{ninna} mean?
\end{quote}

\subsection*{You'll want to know: \te{ఎప్పుడు} \te{వస్తారు}?}
\textbf{\te{ఎప్పుడు} \te{వస్తారు}?} (\emph{eppuḍu vastāru?}) — \textquotedblleft{}when will you come?\textquotedblright{}. The new word in it is \textbf{\te{ఎప్పుడు}}, \textquotedblleft{}when?\textquotedblright{}.

The last gap in the \textbf{\te{ఎ}-} question row: \te{ఎక్కడ} where, \te{ఎవరు} who, \te{ఎందుకు} why, and now \te{ఎప్పుడు} when. The row is complete.

The three days you just learned are its answers, and \textbf{\te{మీరు}} in front of it makes the question respectful.

\begin{grammarlens}[title={What you've built}]
The question the three days answer.
\end{grammarlens}

\subsection*{Your turn}
\begin{itemize}
\raggedright
  \item \emph{Say it:} \emph{eppuḍu}
  \item \emph{Say it:} \emph{eppuḍu}, once more
  \item \emph{Say it:} ask \emph{eppuḍu}, then answer it with \emph{ivāḷa}
  \item \emph{Recall:} say \emph{rātri}, then say \emph{eppuḍu vastāru?}
\end{itemize}

\begin{tcolorbox}[breakable,colback=teal!4,colframe=teal!35!black,title={Before you move on}]
What does \te{ఎప్పుడు} \te{వస్తారు}? mean? (\textquotedblleft{}when will you come?\textquotedblright{}.) Say it once more.
\end{tcolorbox}
\section[mundu]{\te{ముందు} (mundu) — before, in front}

\label{lesson:TE-C75-before}



\begin{quote}
Before the new one: what did \emph{eppuḍu} mean?
\end{quote}

\subsection*{You'll want to know: \te{ముందు}}
\textbf{\te{ముందు}} (\emph{mundu}) — \textquotedblleft{}before\textquotedblright{}.

You already have \te{తరువాత}, \textquotedblleft{}after that\textquotedblright{}. \te{ముందు} is its other half, and until now you could only order events one way.

Like most Telugu relation words it follows its noun: \textbf{\te{భోజనం} \te{ముందు}} — \textquotedblleft{}before the meal\textquotedblright{}. It also means \textquotedblleft{}in front\textquotedblright{} in space, which is the same idea pointed at a room instead of a day.

\begin{grammarlens}[title={What you've built}]
Both ends of an order: \te{ముందు} and \te{తరువాత}.
\end{grammarlens}

\subsection*{Your turn}
\begin{itemize}
\raggedright
  \item \emph{Say it:} \emph{mundu}
  \item \emph{Say it:} \emph{mundu}, once more
  \item \emph{Say it:} say \emph{mundu}, then \emph{taruvāta}, and put a meal between them
  \item \emph{Recall:} say \emph{mariyu}, then \emph{udayam}, then say \emph{mundu}
\end{itemize}

\begin{tcolorbox}[breakable,colback=teal!4,colframe=teal!35!black,title={Before you move on}]
What does \te{ముందు} mean? (\textquotedblleft{}before, in front\textquotedblright{}.) Say it once more.
\end{tcolorbox}
\section[saṁvatsaraṁ]{\te{సంవత్సరం} (saṁvatsaraṁ) — a year}

\label{lesson:TE-C75-year}



\begin{quote}
Before the new one: what did \emph{mundu} mean?
\end{quote}

\subsection*{You'll want to know: \te{సంవత్సరం}}
\textbf{\te{సంవత్సరం}} (\emph{saṁvatsaraṁ}) — \textquotedblleft{}a year\textquotedblright{}.

A Sanskrit tatsama, wearing the same \textbf{-\te{ం}} ending you have seen on \te{ఉదయం} and \te{భోజనం}. Telugu also says \textbf{\te{ఏడు}} for a year of age, but \te{సంవత్సరం} is the plain calendar word.

You now have the hour (\te{గంట}), the day (\te{రోజు}), the month (\te{నెల}) and the year — the whole ladder of Telugu time.

\begin{grammarlens}[title={What you've built}]
The largest unit of time you can name.
\end{grammarlens}

\subsection*{Your turn}
\begin{itemize}
\raggedright
  \item \emph{Say it:} \emph{saṁvatsaraṁ}
  \item \emph{Say it:} \emph{saṁvatsaraṁ}, once more
  \item \emph{Say it:} say \emph{ganṭa}, then \emph{rōju}, then \emph{saṁvatsaraṁ}, smallest to largest
  \item \emph{Recall:} say \emph{kānī}, then \emph{sāyantram}, then say \emph{saṁvatsaraṁ}
\end{itemize}

\begin{tcolorbox}[breakable,colback=teal!4,colframe=teal!35!black,title={Before you move on}]
What does \te{సంవత్సరం} mean? (\textquotedblleft{}a year\textquotedblright{}.) Say it once more.
\end{tcolorbox}
